\section{The Internally Divided Firm and Capacity Formation}
\label{sec:2:divided_firm}

Capacity utilization is not a technical parameter. It is an
inter-scalar political-economic settlement. The macroeconomic
capacity ceiling---the denominator of the utilization rate---is
not a neutral engineering constraint but a historically contingent
outcome conditioned by the balance of class power. Yet identifying
this distributional conditioning at the macro level leaves the
underlying transmission mechanism black-boxed: how, precisely,
does distributive conflict physically alter the conversion of
accumulated capital into productive capacity?

Answering this question requires dismantling a hidden premise
shared by virtually all heterodox growth models: the unified
firm. Whether treating the normal rate of capacity utilization
as an endogenous behavioral target or an exogenous gravitational
anchor, standard frameworks presume a coherent organization
capable of selecting a utilization benchmark, enforcing it across
the labor process, and sustaining it over time. Formally, the
utilization rate is defined as:
%
\begin{equation}\label{eq:utilization}
	\mu_t \equiv \frac{Y_t}{Y^p_t},
\end{equation}
%
\noindent where $Y_t$ is realized output and $Y^p_t$ is potential
productive capacity. But this ratio conceals a deeper question:
who determines $Y^p_t$, and under what organizational conditions?

The unified-firm premise collapses when confronted with the
social organization of capitalist production. In a decentralized
economy, individual capitals commit fixed investments without
\textit{ex-ante} coordination, confronting one another in
competitive markets. Within the enterprise, production decisions
encounter contradictory managerial and worker imperatives.
Capacity utilization cannot be reduced to a target selected by
a single optimizing entity. To theorize the political economy
of productive capacity, I replace the unified firm with two
interconnected frameworks: Okishio's theory of decentralized
market anarchy, and Vidal's theory of the internally divided
firm.

\subsection{Smooth Reproduction and the Unified-Firm Premise}

Proportional reproduction is a condition of feasibility, not a
tendency of capitalist accumulation. Smooth reproduction requires
strict material proportionality between departments of
production---a condition \citet{Basu2022} and \citet{Okishio2022} formally established defined by the latter as the equilibrium accumulation trajectory. \citet{Bergmann2009} call
this balanced trajectory, where realized output meets productive
capacity ($\mu_t = 1$), the \textit{``golden path.''} But the golden path
is a methodological closure, not a gravitational attractor. It
identifies the conditions under which realization can reproduce
itself smoothly; it does not describe the path capitalist accumulation
actually follows. Neither the post-Keynesian assumption of long-run balanced
growth \citep{BleckerSetterfield2019} nor the \textit{``Keynesian in the short run, classical in the long run''} synthesis \citep{DumenilLevy1999} captures this persistence of unbalanced expansion.

Nothing in the logic of decentralized competition
establishes that private investment decisions will reproduce
these proportions. A stylized illustration makes the point
concrete. If productive capacity grows at $5\%$ per year while
aggregate demand for output grows at $2\%$, capacity outpaces
demand by three percentage points each year. After a decade,
utilization falls to roughly $75\%$ of capacity---not because
of a temporary shock, but because uncoordinated investment
persistently builds capacity faster than demand can absorb it.
Balanced growth ($g_Y = g_K$) is merely one configuration of
the transformation elasticity $\theta$; it is not a restoring
force. Smooth reproduction ($\mu_t = 1$) does not require
$\theta = 1$; it holds for any transformation elasticity.
Balanced growth is the narrower special case in which
$\theta = 1$ and $\mu_t = 1$ coincide in the long run.

Because private investment decisions are made without
\textit{ex-ante} coordination, real capitalist economies
routinely violate the proportional conditions of the golden path.
\citet{Okishio2022} demonstrates that chronic disproportionality
is not a temporary friction but the structural norm of
decentralized accumulation. Treating the transient benchmark of
$\mu_t = 1$ as a behavioral target therefore relies on a hidden
premise: that the firm is a unified planning subject capable of
selecting a utilization norm, enforcing it across the labor
process, and sustaining it over time. This premise---the unified
firm---is the target of the present section.

The denominator of the utilization ratio (productive capacity
$Y^p_t$) is not an engineering constant. It is a contested
settlement of divided management and shop-floor struggle. The
numerator (realized output $Y_t$) is hostage to the anarchy of
demand. Smooth reproduction supplies a reference case against
which this historical non-coincidence can be studied. It does not
supply the mechanism that produces it.

\subsection{Decentralized Accumulation and Market Anarchy}

If the ``golden path'' of smooth reproduction is merely a
benchmark of feasibility, the actual trajectory of capitalist
accumulation is governed by the decentralized organization of
investment. Capitalist production is coordinated through
exchange, meaning that the social validation of private
decisions occurs only \textit{ex-post}. Because money mediates
exchange, sale and purchase are separated moments of the capital
circuit \citep{Basu2022}. Surplus value must be realized in
circulation, not by construction \citep{Foley1985, Trigg2020}.
This constitutes the logic of capital as impersonal domination:
capitalists commit to production and investment with the law of
value operating \textit{``behind their backs.''}\footnote{Karl
	Marx, \textit{Capital}, Volume~I, Chapter~1, Section~4 (\textit{``The
	Fetishism of Commodities and the Secret thereof''}), where the
	law of value operates as an overriding law of nature when
	private producers exchange products \citep{MarxKVI}.} They are subject to market
forces and value relations they neither control nor consciously
direct. Critically, this anarchy is not a coordination failure
in an otherwise functional market. The market is a socially
constituted institution whose normal operation produces
disproportionality as a structural feature
\citep{VidalPeck2012}. \citet{Okishio2022} locates the source
of macroeconomic instability precisely in this constitutive
anarchy. Individual capitals commit to fixed investment on the
basis of private profitability expectations, yet the market validates
those commitments only after production has begun. No firm can
know in advance whether its expansion fits the investment and
consumption decisions taken elsewhere. The coordination problem
is not one of imperfect information surrounding an otherwise
coherent plan; no economy-wide plan exists.

Smooth reproduction requires proportional expansion across
sectors producing means of production and sectors producing
consumption goods. Because private investment does not enforce
this proportionality, firms extrapolate past prices, profit
rates, and sales. A signal that induces expansion in one branch
routinely produces excess capacity before complementary demand
materializes elsewhere. Contraction works through the same
cumulative mechanism: falling sales cut investment, reducing
aggregate demand and validating further cuts. Removing Say's
Law reintroduces realization as a substantive condition of
reproduction: accumulation may generate capacity that demand
fails to absorb, or demand may temporarily exceed productive
capacity \citep{Foley1985}. Prices register the
resulting imbalance, but they do not prevent the sunk
commitments produced by it.


Fixed capital gives this temporal lag its structural force.
Structures and machinery embody long-horizon commitments that
cannot be reversed when relative prices or demand conditions
shift. Once installed, excess capacity leaves the firm to
operate below its anticipated scale, defer maintenance, hold
equipment idle, or devalue the asset. These responses
distribute the costs of market anarchy among firms, workers,
and creditors, but they do not restore the missing
\textit{ex-ante} coordination.

Idle capacity is therefore a recurrent, structural form of
decentralized accumulation rather than a temporary friction.
Private investment physically builds capacity before the market
can realize its social use, recurrently a gap
between installed productive capacities ($Y^p_t$) and effective demand ($Y_t$). The capacity utilization rate ($\mu_t$) registers this gap at the macro level. Unbalanced growth is not a disturbance around proportionality; it is the structural consequence of distributively conditioned investment under
decentralized decision-making. The analytical problem is
therefore not how proportionality re-emerges, but how persistent
imbalance remains historically reproducible. The condition
$\mu_t = 1$ is at most a transient point through a given trajectory
proportionality, while persistent divergence records the ongoing
reality of market anarchy.

Okishio's macro framework stops one step short: it does not explain how firms translate this external anarchy into internal operating decisions. The firm does not manage market anarchy; it internalizes it. A firm does not receive a price signal and respond as a frictionless, unified unit. It processes the signal through an organization divided over control, cooperation,
effort, and the allocation of resources. Vidal's internally
divided firm provides that missing mediation.

\subsection{The Internally Divided Firm}


Vidal builds on Marx's Volume~I theory of the labor process, drawing on contemporary organizational theory and sociological institutionalism to develop an organizational political economy of the capitalist firm \citep{VidalPeck2012, Vidal2022}. Marx established that
capitalist management faces a constitutive contradiction: it
must simultaneously organize the cooperative labor required for
production and enforce the coercive control required for
valorization. \citet{Vidal2022} demonstrates how this contradiction fractures
the modern capitalist firm, treating it as a divided social
organization whose internal structure is shaped by cognitive,
normative, regulative, and organizational institutions. Management must extract surplus value through
control over labor, but production also requires workers'
cooperation, tacit knowledge, and problem-solving capacity.
Standardization and surveillance secure a measurable pace of
work while weakening the discretion needed to sustain complex
production. The firm cannot maximize both imperatives at once
because the labor process that produces value is also the site
where workers contest its extraction. This contradiction is not
a universal constant of capitalism; it is institutionally
mediated, taking variegated forms across distinct accumulation regimes throughout historical capitalism
\citep{Vidal2019}.

Workers respond to this contradiction by actively calibrating
effort, cooperation, machine pacing, and the deployment of tacit
knowledge. Capitalist management extracts potential capacity
through concrete shop-floor mechanisms: increasing machine speed,
deferring maintenance to hit short-term throughput targets, and
enforcing rigid surveillance. Workers contest this extraction
through labor turnover, collective exit, or the deliberate
withdrawal of problem-solving cooperation and effort. Throughput,
quality, and maintenance therefore depend on social relations
that no engineering specification fixes in advance. A production
target that raises machine speed by deferring maintenance may
increase current output while degrading sustainable capacity. A
directive that suppresses worker discretion may meet a
short-term labor-control objective while worsening workflow
bottlenecks. Labor effort is consequently not a fixed technical
parameter applied to a pre-existing capital stock; it is an
active, contested determinant of the capacity utilization itself.

This dynamic enforces a strict wedge between
\textit{engineering capacity} and \textit{practical capacity}.
Engineering capacity records the theoretical maximum output of
equipment under ideal technical conditions. Practical capacity
records the sustainable output rate dictated by actual
maintenance schedules, product-mix complexity, and the negotiated
intensity of labor. When managers report capacity to national
statistical agencies, they rely on conventions shaped by their
organizational position and the immediate pressures of the shop
floor. Aggregate capacity statistics therefore do not reveal a
neutral physical ceiling; they aggregate situated, institutional
judgments \citep{Shaikh2016, Nikiforos2016}. Complementarily,
empirical surveys of managerial practice confirm that firms do
not systematically organize production decisions around a
coherent capacity utilization target \citep{Singh2022}. For the
purposes of macroeconomic modeling, the productive capacity
ceiling ($Y^p_t$) is not an exogenous engineering constraint but
an endogenous settlement of workplace conflict. This settlement
is the micro-foundation through which class struggle and
distribution transmit into capacity utilization, setting up the
distributive--technological mapping developed in the following
subsection.

Because the valorization--labor process contradiction cannot be
solved, firms settle into organizational configurations of
\textit{``mediocre sufficiency''} \citep{Vidal2022}. These are hybrid
arrangements that reproduce production without resolving the
underlying conflict between control and cooperation. Partial
worker autonomy coexists with rigid discipline; continuous
improvement programs coexist with brutal throughput pressure;
machinery might remain idle alongside binding productive bottlenecks. Following
\citet{Vidal2014}, these arrangements are best understood as
\textit{institutional incoherence}: the firm's internal
settlements are functional enough to sustain accumulation yet
contradictory enough to generate persistent friction. This stands
in direct contrast to the Varieties of Capitalism assumption of
institutional complementarity, where national institutions are
presumed to form coherent, mutually reinforcing systems \citep{Hall2003}. Here,
incoherence is not a deviation from an optimizing norm; it is the
permanent, structural form through which capitalist management
physically converts accumulation demand into capacity usage. 

\subsection{Distribution, Technique, and the Transformation
	of Accumulation}

The choice of technique connects the divided firm to
distribution, governing how accumulated capital is physically
transformed into productive capacity. Capitalists adopt a new
technique only when it reduces the cost of production at
prevailing prices and wages; such a technique is capitalistically
viable \citep{Okishio1961}. \citet{Okishio1961} demonstrates
that a viable capital-using, labor-saving (CU-LS) technique
raises the equilibrium profit rate when the real wage remains
constant---a direct challenge to Marx's law of the tendential
fall in the rate of profit. Marx's account, by contrast, allows
the same form of technical change to depress profitability when
rising capital intensity increases the capital advanced relative
to surplus value. \citet{Basu2022} reconciles this seemingly mutually exclusive 
theorizations, by formalizing the \textit{Marx--Okishio threshold}: 
the rate of real-wage growth defines the result driven by
the balance of class forces. 

This threshold situates the present framework within the
induced-innovation tradition while departing from it. The
Innovation Possibility Frontier (IPF) literature, initiated by
\citet{Kennedy1964}, treats distribution as selecting the
\textit{direction} of technical change along a given frontier:
relative factor costs determine whether innovation is
labor-augmenting or capital-augmenting, while the frontier
itself---the scale of feasible productivity growth---remains
exogenous \citep{Julius2005, Tavani2012}. Later assessments
sharpen this limitation: IPF-based models explain the
composition of technical change but remain largely silent on
innovation intensity \citep{TavaniZamparelli2018}.
\citet{Zamparelli2024a} partially endogenizes scale by mapping
a family of IPF levels to R\&D effort, yet the analytical
structure continues to rely on frontier geometry and
steady-growth closure \citep{Zamparelli2024b}. An alternative
strand reframes inducement through mechanization rather than
frontier selection: \citet{Kaldor1957} links productivity
growth to capital deepening, while
\citet{CajasGuajirro2024} reformulates the IPF with an explicit
choice over mechanization growth under distributive and
labor-market conditions. The present framework builds upon a 
mechanization productive frontier,  treating the distributive--technological 
relation as a regime-specific transformation between accumulation and
productive capacity. Rather than selecting along a given
frontier, distribution conditions mechanization intensity,
which in turn determines the transformation elasticity linking
accumulation to the formation of productive capacities. 

The Marx--Okishio threshold gives wage pressure a precise role
in mechanization, but the profitability outcome depends on the
labor market closure. \citet{Basu2026} demonstrates that the
equilibrium profit rate after viable CU-LS technical change is
governed by the institutional structure of capital--labor
bargaining. Under a constant-wage-share closure---what
\citet{Basu2026} formalizes as the ``strong labor'' regime,
applicable to Atlantic Fordism---any viable CU-LS technique
unambiguously \textit{lowers} the equilibrium profit rate.
Marx's LTFRP reigns supreme precisely where labor is powerful
enough to compel real wages to grow in tandem with productivity.
Under this closure, rapid mechanization is induced by
distributive pressure, yet the same pressure contests capital's
claim on the resulting productivity gains
\citep{Vidal2015, Vidal2019}. A stylized illustration makes the
mechanism concrete. If a $10\%$ real wage increase raises
mechanization intensity by $8\%$ under strong-labor conditions,
the transformation elasticity $\theta$ rises from $0.70$ to
$0.85$: accumulation translates more efficiently into productive
capacity because firms substitute machinery for labor. Wage-led
mechanization was not a harmonious productivity bargain; but a
a contested settlement in which the distributive pressure that
induced technical change simultaneously eroded the profit rate
it was meant to restore.

The weakening of organized labor after the Fordist crisis
altered that mechanism. Under a \textit{``weak labor''} closure---where
the unemployment rate remains constant such that workers lack the
institutional power to enforce a target wage share---the
distributive pressure to replace labor is suppressed
\citep{Basu2026}. Yet mechanization does not necessarily cease. Lower
relative machinery prices and technical change in capital-goods
production can continue to favor capital-using techniques through a
second, independent channel. Distribution and machinery prices
therefore constitute a dual-channel mechanism in the
choice of technique. Their relative force changes across
accumulation regimes, which prevents the wage share from being
treated as an external control added to an otherwise autonomous
technical relation. Extending the stylized example: under
Post-Fordist conditions, a $15\%$ decline in relative machinery
prices raises mechanization intensity by $6\%$, partially
compensating for the weakened distributive channel and
sustaining $\theta$ near $0.75$ despite wage suppression. The
same $10\%$ wage increase that produced an $8\%$ mechanization
response under Fordism produces only a $2\%$ response under
Post-Fordism, leaving $\theta$ near $0.72$. The two channels do
not simply add; they interact through the institutional
configuration that governs their relative force.

These channels still operate through the contradictions
identified in the preceding subsections. Firms choose techniques
before demand validates the associated capacity, and they
operate those techniques through labor processes that divide
control from productive cooperation. Investment consequently
does not translate one-for-one into a stable productive ceiling.
The transformation depends on the composition of capital, the
distributive conditions governing mechanization, and the
organization through which installed equipment becomes usable
capacity. In a labor surplus economy---the third closure
identified by \citet{Basu2026}, applicable to peripheral
capitalism---the real wage remains constant and the equilibrium
profit rate rises unambiguously after viable CU-LS technical
change. Yet the peripheral transformation is structurally
truncated by external financial subordination and the balance of
payments constraint, a mechanism developed in the following
subsection.

Okishio and Vidal jointly rule out a direct
passage from accumulation to productive capacity.
Decentralized investment generates structural
disproportionality across firms and sectors; the internally
divided firm internalizes this disproportionality as
contested operating conditions rather than a coherent
utilization target. The analytical object that captures this
mediated transformation is the \textit{transformation elasticity} ($\theta \equiv d y^p / d k$), which governs how
accumulation converts into productive capacity under a context of
distributive conflict and market anarchy. The empirical
analysis limits the original inter-departamental formulation by Marx in Capital Volume II into a
single-sector aggregate, masking internal sectoral
divergences while preserving the core logic.

\subsection{Historical Regimes of Institutional Mediation and the Transformation Elasticity}

The transformation elasticity is an emergent property of class struggle 
that mediates the process of capital accumulation and its transformation 
into productive capacities under the prevailing institutional environment.
Decentralized accumulation does not automatically guarantee
proportional reproduction ($\theta = 1$). Institutional
mediation---wage-setting norms, financial architectures, and
competitive structures---conditions how distributive pressures
translate into mechanization choices
\citep{Aglietta2000}. 

\citet{Basu2026} demonstrates that the equilibrium profit
rate after viable capital-using, labor-saving (CU-LS)
technical change is governed by the labor market closure,
which maps directly onto distinct historical accumulation
regimes. Under a \textit{constant-wage-share} closure---the
institutional signature of Atlantic Fordism
(1947--1973)---workers possess the bargaining power to grasp
productivity gains, forcing real wages to grow in tandem with
labor productivity. In this strong-labor regime, viable
mechanization unambiguously lowers the equilibrium profit
rate, so that Marx's law of the tendential fall in the rate
of profit holds \citep{Basu2026}. This stagnation tendency
was institutionally offset: the class compromise, mass
consumption, oligopolistic market structures, and military
Keynesianism of the Fordist era absorbed the surplus and
sustained effective demand \citep{Vidal2019}.

Under a \textit{constant-unemployment} closure---which resembles the persistent unemployment feature of the post-Fordist accumulation regime---an era during which the class compromise that took place during the Fordist period is dismantled and labor's
bargaining power is severely weakened. The wage share falls,
and mechanization is driven not by distributive pressure but
by the relative cheapening of ICT and the shift toward
financially reversible, short-horizon assets, patents and intellectual property of intangible assets 
\citep{Lapavitsas2011, Pagano2014}. Because wages no longer
track productivity, the profit rate may stabilize or rise,
but the suppression of wage-led demand generates a chronic
under-consumptionist stagnation tendency. Following \citet{Vidal2019}, I hypothesize that under structural conditions, the post-fordist accumulation regime produces a dysfunctional utilization configuration:
accumulation decouples from productive capacity formation,
generating chronic idle capacity that is temporarily masked
by debt-led consumption and financial fragility rather than
sustainable realization \citep{Stockhammer2012, Vidal2019}.

In the capitalist periphery, different closures represent distinct historical configurations throughout Chilean history. First, under an assumption of a
\textit{constant-real-wage}, which corresponds to
Okishio's original formulation maps onto a surplus-labor
economy \citep{Lewis1954, Basu2026}. During the
import-substituting industrialization (ISI) era, the
\textit{rural-urban} and \textit{demographic} transition enabled by the transformation of sanitary conditions and a massive drop of child mortality in the early XX century,  regulated late industrialization, providing an unlimited supply of labor from the countryside that disciplined urban wages and allowed the profit rate to reproduce itself alongside mechanization. 
However, peripheral mechanization is structurally import-biased:
widening the mechanization gap requires imported capital
goods and technologically non-substitutable inputs
\citep{Prebisch1981}. The balance of payments constraint
therefore acts as a hard external ceiling on the
transformation of accumulation into capacity. Under a
developmentalist regime with surplus labor (1933--1957), this operated as an
\textit{import capacity ceiling}. Second, a \textit{constant-wage-share} or more accurately \textit{rising-wage-share} closure, represents the intensification of stagnation tendencies leading to structural crisis of capitalism once the unlimited supply of labor depleted and the incentives for capitalist firms to mechanize risen due distributive pressures while increasing import requirements and stressing further the external constraint. Third, the discipline imposed by the neoliberal counter-revolution \citep{Casals2021} (1974--present) a \textit{constant-unemployment} closure coupled with trade and capital account liberalization relieved the external constraint, but locked the periphery into extractive specialization and premature deindustrialization, such that the WTO and bilateral free
trade agreements legally encoded a \textit{policy space ceiling} \citep{Ahumada2019, Palma2019}.


These historical configurations establish that the productive
ceiling is not a neutral, trans-historical engineering
constraint but a contingent outcome of decentralized
accumulation, class power, and a country's relative
positionality within the capitalist world system. The
center--periphery comparison reveals a fundamental asymmetry:
in the center, structural crises arise from the
\textit{exhaustion} of institutional mediation (e.g., the
1970s profit squeeze); in the periphery, arise from
the \textit{absence} of sovereign mediation and the
truncation imposed by the hierarchically oranized financial order imposing 
an external constraint that is not necessarily binding. 